\begin{solapartat}
\begin{minipage}[t]{0.72\linewidth}
\punts{0,3} $\Delta V = \dfrac{E_C}{e} = \dfrac{1000\ \mathrm{eV}}{e} = 1000\ \mathrm{V}$

\punts{0,4} $\left|\vec{E}\right| = \dfrac{\Delta V}{d} = 1000\ \mathrm{V/m} = 1000\ \mathrm{N/C}$

\punts{0,3} Els electrons parteixen del repòs, de manera que descriuen una trajectòria rectilínia i la força té la mateixa direcció i sentit que el desplaçament de la càrrega. Com que la càrrega de l'electró és negativa, el sentit d'$\vec{E}$ és l'oposat del de la força $\vec{F}$.
\end{minipage}\hfill
\begin{minipage}[t]{0.24\linewidth}
\vspace{0pt}
\figura[1]{sol_fig1.png}
\end{minipage}
\end{solapartat}

\begin{solapartat}
\punts{0,3} $E_C = 1000\ \mathrm{eV} = 1,602\times 10^{-16}\ \mathrm{J}$
\[ E_C = \frac{1}{2}mv^2 \Longrightarrow v = \sqrt{\frac{2\cdot 1,602\times 10^{-16}}{9,11\times 10^{-31}}} = 1,875\times 10^{7}\ \mathrm{m/s} \]

\punts{0,4} $\left|\vec{F}\right| = e\cdot v\cdot B\cdot \sin\left(\dfrac{\pi}{2}\right) = 4,51\times 10^{-13}\ \mathrm{N}$

\begin{minipage}[t]{0.78\linewidth}
\punts{0,3} És necessari que en la justificació o es faci un esquema de les forces i de la trajectòria o s'indiqui que la força magnètica és perpendicular a la velocitat, per tant, només hi ha la component normal de l'acceleració. A més, aquesta acceleració és de mòdul constant, per tant, la trajectòria serà circular.
\end{minipage}\hfill
\begin{minipage}[t]{0.18\linewidth}
\vspace{0pt}
\figura[1]{sol_fig2.png}
\end{minipage}
\end{solapartat}
